\chapter{Implementation}
\ChapterLabel{Implementation}
\section{Chapter Responsibility}
\label{Section:ImplementationResponsibility}
Make the architecture reproducible through concrete implementation guidance.

\noindent\textbf{Expected result:} A repeatable path from design to a working implementation.

\section{Revision 0.1 Scope Notice}
\label{Section:DeferredImplementation}
\begin{fsaNote}[Deferred project content --- template reference only]
This chapter is retained in the complete document as template guidance.
Its sample artifacts, IDs, priorities, statuses, versions, outcomes, and
approval examples do not describe the FTDI245R project. In particular, any
sample Done, Pass, or release entry is not implementation, verification,
validation, or approval evidence for this project. Detailed project content
is deferred under \hyperref[Section:ProjectScope]{the revision scope}.
See \hyperref[Table:ProjectGapReview]{the reviewed gap list} for actual
decisions and unresolved items.
\end{fsaNote}

\section{How to Use This Chapter}
Document the selected technology baseline, build and run instructions, coding conventions that enforce architecture, configuration, and definition of done. Keep design rationale in Architecture Overview.
\section{Example}
The implementation uses [language], [framework], and [database]. Build with \texttt{[command]}; run with \texttt{[command]}. Configuration comes from [approved source]; credentials are never stored in source control. Completion requires passing tests, reviewed traceability, and approved evidence.
